% Part: computability
% Chapter: computability-theory
% Section: reducibility

\documentclass[../../../include/open-logic-section]{subfiles}

\begin{document}

\olfileid{cmp}{thy}{red}
\olsection{হ্রাসযোগ্যতা}

\begin{explain}
এখন আমরা জানি, অন্তত একটি সেট~$K_0$ গণনসাধ্যভাবে তালিকায়নযোগ্য, কিন্তু
গণনসাধ্য নয়। এমন আরও সেট আছে—এটি স্পষ্ট হওয়া উচিত। হ্রাসযোগ্যতার
পদ্ধতি বারবার প্রথম নীতি থেকে শুরু না করেই অন্য সেটের এই বৈশিষ্ট্যগুলি
প্রমাণ করার একটি শক্তিশালী উপায় দেয়।

সাধারণভাবে, কোনো সেট~$A$-কে সেট~$B$-তে “হ্রাসকরণ” হলো, !!{element}s
$B$-তে আছে কি না সেই প্রশ্নের উত্তরকে !!{element}s $A$-তে আছে কি
না সেই প্রশ্নের উত্তরে রূপান্তর করার একটি পদ্ধতি। আমরা “বহু-এক
হ্রাসযোগ্যতা” নামের একটি ধারণায় মন দেব; তবে ভিন্ন ভিন্ন বৈশিষ্ট্যসহ
হ্রাসযোগ্যতার আরও অনেক ধারণা আছে। গণনামূলক জটিলতা অধ্যয়নেও হ্রাসযোগ্যতার
ধারণা কেন্দ্রীয়, যেখানে দক্ষতার প্রশ্নও বিবেচনা করতে হয়। যেমন, কোনো
সেটকে “NP-সম্পূর্ণ” বলা হয় যদি সেটি NP-তে থাকে এবং প্রত্যেক NP সমস্যাকে
তার মধ্যে হ্রাস করা যায়; এখানে নিচে বর্ণিত হ্রাসকরণের অনুরূপ একটি ধারণা
ব্যবহৃত হয়, শুধু অতিরিক্ত শর্ত থাকে যে হ্রাসকরণটি বহুপদী সময়ে গণনা করা
যাবে।

হ্রাসকরণের ধারণাটি আমরা ইতিমধ্যেই পরোক্ষভাবে ব্যবহার করেছি। সেট~$K$-কে
সংজ্ঞায়িত করি:
\[
K = \Setabs{x}{\cfind{x}(x) \fdefined},
\]
অর্থাৎ, $K = \Setabs{x}{x \in W_x}$। থামার সমস্যা অসমাধানযোগ্য—এই
বিষয়ে আমাদের প্রমাণটি (\olref[hlt]{thm:halting-problem}) সবচেয়ে সরাসরি
দেখায় যে~$K$ গণনসাধ্য নয়। মনে করুন, $K_0$ হলো সেট
\[
K_0 = \Setabs{\tuple{e, x}}{\cfind{e}(x) \fdefined },
\]
অর্থাৎ, $K_0 = \Setabs{\tuple{e,x}}{x \in W_e}$। $K$-এর
অগণনসাধ্যতার যেকোনো প্রমাণকে সহজেই~$K_0$-এর অগণনসাধ্যতার প্রমাণে
বাড়ানো যায়: $K_0$ গণনসাধ্য হলে, !!a{element}~$x$, $K$-তে আছে কি না
আমরা শুধু জিজ্ঞাসা করেই নির্ণয় করতে পারতাম যে জোড়া~$\tuple{x, x}$,
$K_0$-তে আছে কি না। $x$-কে~$\tuple{x, x}$-এ পাঠায় এমন অপেক্ষক~$f$,
$K$ থেকে~$K_0$-তে একটি \emph{হ্রাসকরণের} উদাহরণ।
% BN-SRC-152 TARGET-CORRECTION-BEGIN
\par\smallskip\noindent\textnormal{\small\textbf{উৎস-সংশোধন (BN-SRC-152):} হিমায়িত ইংরেজি পাঠে $K_0$-এর সমতুল্য দ্বিতীয় সেট-গঠনে জোড়াটি $\tuple{x,e}$ লেখা হয়েছে, অথচ প্রথম সংজ্ঞা ও শর্ত $x\in W_e$ অনুসারে সূচকটি প্রথম এবং নিবেশটি দ্বিতীয় উপাদান হওয়া দরকার। বাংলা পাঠে $\tuple{e,x}$ রাখা হয়েছে।} % BN-SRC-152 TARGET-CORRECTION-END
\end{explain}

\begin{defn}
$A$ ও $B$-কে স্বাভাবিক সংখ্যার সেট ধরি। একটি গণনসাধ্য
অপেক্ষক~$f\colon \Nat \to \Nat$, $A$ থেকে~$B$-তে একটি \emph{বহু-এক
হ্রাসকরণ} যদি এবং কেবল যদি প্রত্যেক স্বাভাবিক সংখ্যা~$x$-এর জন্য
\[
x \in A \quad \text{যদি এবং কেবল যদি} \quad f(x) \in B.
\]
এমন হ্রাসকরণ~$f$ থাকলে বলি, $A$, $B$-তে \emph{বহু-এক হ্রাসযোগ্য}; লিখি
$A \leq_m B$। $A$, $B$-তে এবং বিপরীতভাবেও বহু-এক হ্রাসযোগ্য হলে $A$
ও~$B$-কে \emph{বহু-এক সমতুল্য} বলা হয়; লিখি $A \equiv_m B$।
\end{defn}

উপরের সংজ্ঞার অপেক্ষক~$f$ যদি এক-এক হয়, তবে $A$-কে~$B$-তে
\emph{এক-এক হ্রাসযোগ্য} বলা হয়। নিচে বর্ণিত অধিকাংশ হ্রাসকরণ এই অধিকতর
শক্তিশালী শর্ত মেটায়, কিন্তু আমরা এই তথ্যটি ব্যবহার করব না।

\begin{digress}
এক-এক হ্রাসযোগ্যতা যে সত্যিই বহু-এক হ্রাসযোগ্যতার চেয়ে শক্তিশালী শর্ত,
তা সত্য হলেও মোটেই স্পষ্ট নয়। অন্যভাবে বললে, এমন অসীম সেট $A$ ও~$B$
আছে, যেখানে $A$, $B$-তে বহু-এক হ্রাসযোগ্য, কিন্তু $B$-তে এক-এক
হ্রাসযোগ্য নয়।
\end{digress}

\end{document}
